% !TeX program = xelatex
% ============================================================
% interim_report/interim_report.tex — «Отчёт ВКР» — деливерабл КТ2
% исследовательской (академической) ВКР ОП «Программная инженерия».
%
% Жанр задан «Инструкцией по подготовке ВКР КТ2 2025/26»: документ по
% ГОСТ 7.32-2017, НЕ МЕНЕЕ 15–20 страниц со списком источников (>= 5 по
% ГОСТ). Несёт «проработанный Подход» к исследованию. Разделы — ровно
% перечень инструкции; если по какой-то части сведений нет — она
% опускается, ПОСЛЕДОВАТЕЛЬНОСТЬ остальных сохраняется. Титульный лист
% подписывают студент и руководитель (раздел «Подписи»).
%
% ЭТО НЕ ТЕКСТ ВКР (тот — документ thesis: реферат, главы, заключение).
% Отчёт ВКР — развитие Аннотации КТ1: те же части + экспериментальная
% база, описание данных и подход к экспериментам.
%
% Как пользоваться: замените жёлтые \hseFill{…} своим текстом и
% разворачивайте каждый раздел до содержательного объёма (суммарно
% 15–20 стр.). Данные о себе, руководителе, УДК и годе — в настройках.
% ============================================================
\documentclass[12pt,a4paper]{extarticle}
% !TeX program = xelatex
% ============================================================
% common/annotation_preamble.tex — преамбула Аннотации (документ КТ1
% исследовательской ВКР ОП «Программная инженерия»).
%
% ЭТО НЕ ЕСПД-документ: без штампа/рамки, без кода документа, без листа
% утверждения и без содержания. Оформление по ГОСТ 7.32-2017: поля
% 30/15/20/20, Times New Roman 12 пт, интервал 1,5, обычный номер
% страницы снизу. Данные — из common/meta (\hse-макросы). Ссылки [n]
% печатаются В СТРОКУ, список источников — по ГОСТ Р 7.0.5-2008.
% ============================================================

% ── Шрифты (автогенерируются из настроек проекта) ───────────
\newcommand{\hseDefaultMono}{Courier New}
\input{../common/fonts}
\usepackage[english,russian]{babel}

% ── Геометрия по ГОСТ 7.32-2017 §6.1.1 (без места под штамп) ─
\usepackage[a4paper,left=30mm,right=15mm,top=20mm,bottom=20mm]{geometry}
\usepackage{setspace}
\onehalfspacing
\usepackage{indentfirst}
\setlength{\parindent}{1.25cm}
\setlength{\parskip}{6pt}

% ── Типографика и таблицы ───────────────────────────────────
\usepackage{microtype}
\usepackage{csquotes}
\usepackage{enumitem}
\setlist[itemize,1]{label=--}
\usepackage{graphicx}
\usepackage[table]{xcolor}
\usepackage{array}
\usepackage{tabularx}
\usepackage{booktabs}
\usepackage{float}

% ── URL и переносы ──────────────────────────────────────────
\usepackage{url}
\usepackage{xurl}
\urlstyle{same}
\usepackage{hyphenat}

% ── Цитирование: ссылки [n] В СТРОКУ (как в реальных аннотациях),
%    список источников — ГОСТ Р 7.0.5-2008. НЕ переопределяем \@cite
%    в надстрочный вид (в отличие от ЕСПД-преамбулы). ────────────
\usepackage{cite}
\makeatletter
\renewcommand\@biblabel[1]{#1.}
\makeatother

\usepackage{lastpage}
\usepackage{hyperref}
\hypersetup{
    unicode=true,
    breaklinks=true,
    colorlinks=true,
    linkcolor=black,
    urlcolor=blue,
    citecolor=blue,
    pdfborder={0 0 0}
}

% ── Данные проекта (\hse…), затем «человеческие» имена (commands) ──
% meta даёт \hseFill/\hseOptional и все \hse-макросы; commands —
% \signdateblank, \hseExecutorsBlock-зависимости (\authorsignatureline,
% \studentgroup, \studentname) для блока «Выполнил» в титуле.
\input{../common/meta}
\input{../common/commands}

% Без штампа: обычный номер страницы снизу по центру (ГОСТ 7.32 §6.3.1).
\pagestyle{plain}

\begin{document}
\sloppy
\input{title}


\begin{center}
  \textbf{\large Отчёт ВКР}
\end{center}

\vspace{0.4em}

% ── Ключевые слова (ГОСТ 7.32-2017 §5.3.2.1 / §6.12.2): 5–15 слов или
%    словосочетаний, ПРОПИСНЫМИ, в именительном падеже, через запятую,
%    без точки в конце. ──
\noindent\textbf{Ключевые слова:} \hseFill{5–15 СЛОВ ПРОПИСНЫМИ В ИМЕНИТЕЛЬНОМ ПАДЕЖЕ ЧЕРЕЗ ЗАПЯТУЮ БЕЗ ТОЧКИ В КОНЦЕ}

\section{Оценка современного состояния решаемой проблемы}
Рассматриваемая проблема состоит в~\hseFill{сформулируйте научную проблему}. Сегодня для её решения применяются \hseFill{существующие подходы, методы, работы со ссылками} \cite{example-en}. Их основные ограничения: \hseFill{ограничение 1}; \hseFill{ограничение 2}. \hseFill{Разверните обзор: сравнение подходов, таблицы, выводы — это основной «обзорный» раздел отчёта}.

\section{Основание и исходные данные для исследования}
Основанием для выполнения работы является учебный план ОП «\hseSpecialization» и~тема ВКР, утверждённая приказом ФКН НИУ ВШЭ. Исходные данные: \hseFill{имеющийся задел, данные, публикации, программные средства, от которых отталкивается работа}.

\section{Актуальность и новизна темы исследования}
Актуальность обусловлена \hseFill{практическая и научная значимость задачи}. Научная новизна работы заключается в~\hseFill{что именно предлагается впервые}.

\section{Объект и предмет исследования}
\textbf{Объект исследования (или разработки)}~--- \hseFill{что изучается в целом}.

\textbf{Предмет исследования}~--- \hseFill{конкретная сторона объекта, на которой сфокусирована работа}.

\section{Цель исследования}
\hseFill{одна формулировка достигаемого научного результата}.

\section{Задачи исследования}
\begin{enumerate}[label=\arabic*), leftmargin=1.3cm, topsep=2pt, itemsep=2pt]
  \item \hseFill{провести обзор существующих подходов};
  \item \hseFill{предложить и формализовать метод (модель, алгоритм)};
  \item \hseFill{реализовать программное средство для экспериментов};
  \item \hseFill{провести вычислительный эксперимент и оценить результаты};
  \item \hseFill{сделать выводы о применимости}.
\end{enumerate}

\section{Методы и методология проведения исследования}
\hseFill{опишите применяемые методы и общую методологию: аналитический обзор, формальное моделирование, вычислительный эксперимент, статистический анализ; как они связаны с задачами}.

\section{Ожидаемые новые результаты исследования}
\begin{enumerate}[label=\arabic*), leftmargin=1.3cm, topsep=2pt, itemsep=2pt]
  \item \hseFill{метод / модель / алгоритм и в чём его новизна};
  \item \hseFill{программная реализация для проведения экспериментов};
  \item \hseFill{экспериментальная оценка и сравнение с базовыми методами};
  \item \hseFill{выводы о границах применимости}.
\end{enumerate}

% ── «Проработанный Подход» (ядро КТ2): база, данные, эксперименты ──
\section{Экспериментальная база исследования}
\hseFill{на чём проводится исследование: вычислительная среда, оборудование, программная инфраструктура, стенды}.

\section{Описание данных для исследования}
\hseFill{данные, на которых будет проведено исследование: источник, объём, структура, способ получения и подготовки, ограничения использования}.

\section{Подход к проведению экспериментов и интерпретации результатов}
\hseFill{план экспериментов: какие серии, какие метрики качества, базовые методы для сравнения, как будут интерпретироваться и проверяться результаты (статистическая значимость, воспроизводимость)}.

\section{Область применения результатов исследования}
\hseFill{укажите область применения}.

\section{Предполагаемое внедрение результатов исследования}
\hseFill{где и в каком виде планируется внедрить результаты: продукт, сервис, публикация, открытый код}.

\section{Экономическая эффективность или значимость исследования}
\hseFill{где и кем применимы полученные результаты; в чём практический или экономический выигрыш}.

\section{Прогнозные предположения о развитии объекта исследования}
\hseFill{как может развиваться объект исследования и какое продолжение возможно у работы}.

\bigskip
% ── Список источников по ГОСТ Р 7.0.5-2008 (методичка ФКН/ПИ): единый
%    алфавитный нумерованный список, сначала латиница, затем кириллица;
%    разделитель областей «. –»; для сетевых ресурсов — «[Электронный
%    ресурс] … Режим доступа: URL (дата обращения: ДД.ММ.ГГ)». На каждый
%    источник должна быть ссылка \cite{…} в тексте.
%    Инструкция КТ2: источников — НЕ МЕНЕЕ ПЯТИ. ──
\noindent\textbf{Список использованных источников}
\vspace{-0.3em}
\renewcommand{\refname}{}
\begin{thebibliography}{99}
  \bibitem{example-en}
  Author, A.~B. Title of the paper / A.~B.~Author, C.~D.~Coauthor // Proceedings of the Conference.~--- New York, NY, USA: ACM.~--- 2024.~--- С.~10--20.

  \bibitem{example-web}
  Название ресурса [Электронный ресурс].~--- Режим доступа: \url{https://example.org} (дата обращения: 01.01.2026).
\end{thebibliography}


\end{document}
